\section{Applications to Hodge spectra and HYM dynamics}
\label{sec:integrated-theorem}

We apply the reconstruction theorem to the analytic family of
Proposition~\ref{prop:analytic-reapplication}. The integral map
recovered from periods is the vanishing-cycle map in the exact
sequence underlying the Hodge comparison. We then identify the
spectral, dynamical and quotient constructions of \cite{W_HYM}
in common coordinates: the invariant graph gives the finite Hessian,
and the infinitesimal unitary action gives both the weak section
energy and the vertical directions in the harmonic quotient.

Let $V=H^2(Y_{h_*};\RR)$ and let $G$ be its central harmonic
mass. The subspace $K_H$ of classes surviving on the smoothing
has dimension two; in the fixed divisor coordinates it is
spanned by $k_1=g+\xi_1$ and $k_2=g+\xi_2$.
Represent $V/K_H$ by $W_H=K_H^{\perp_G}$, using the
$G$-orthogonal lifts of $f,g$ as coordinates.
Write $V_{\rm lost}=V/K_H$.
The curve restriction map descends to
\begin{equation}\label{eq:lost-restriction}
 \mathsf r:V_{\rm lost}\longrightarrow\RR^8,\qquad
 \mathsf r(x_f,x_g)=x_fv_f+x_gv_g.
\end{equation}
Here $v_f,v_g$ have the orientations fixed in
\eqref{eq:integral-attachment}.

Put
\begin{equation}\label{eq:hodge-model}
 \kappa_i=\frac{2\pi}{|\log|\tau_i||},\qquad
 w_\alpha=\sum_{i\in I_\alpha}\kappa_i,\qquad
 E_H=\begin{pmatrix}w_f&0\\0&w_g\end{pmatrix},\qquad
 M_H=\frac45\begin{pmatrix}9&-1\\-1&9\end{pmatrix}.
\end{equation}
Here $E_H$ comes from the capacity-weighted link restriction
and $M_H$ from harmonic mass after removing the full surviving
space. Neither is defined from the small eigenvalues.

For the bundle application, set
$\mathcal W=\operatorname{Mat}_2(\CC)^{\oplus2}$ and
$V_H=H^3/6=160/3$.
Fix the two algebraic Koszul extension frames. Their harmonic
representatives have squared masses $c_+(p),c_-(p)>0$,
uniformly bounded above and below and converging to positive
central limits. These are elliptic data for $g_p$ and its HE
line backgrounds; intersection numbers do not compute them.
Let $K_B=U(2)\times U(2)$ act by
$(k_+,k_-)(A,B)=(k_+Ak_-^{-1},k_-Bk_+^{-1})$.
On
\[
 \mathfrak h=\{(S_+,S_-):S_\pm=S_\pm^*,\
                              \tr S_++\tr S_-=0\}
\]
define
\begin{equation}\label{eq:bundle-model}
 \begin{split}
 M_B(S,T)&=V_H\tr(ST),\\
 E_B(p,z)(S)&=2c_+\|S_+A-AS_-\|_{\rm HS}^2
                    +2c_-\|S_-B-BS_+\|_{\rm HS}^2,\\
 \mu_{2,p}(z)&=V_H^{-1}
 (c_+AA^*-c_-B^*B,\ c_-BB^*-c_+A^*A),\qquad z=(A,B).
 \end{split}
\end{equation}

The real equivariant normalization $w=\chi_p(z)$ and
section operator $\mathcal L_{p,w}$ are specified in
Recollection~\ref{thm:bundle-input}. The invariant graph
$Z(S)=S+\psi(S)$ and split projection $\Pi$ are those of
Recollection~\ref{thm:flow-input}.
On the real tangent space of $\mathcal W$ put
\begin{equation}\label{eq:arrow-mass}
 G_{0,p}(q,q)=2c_+(p)\|q_A\|_{\rm HS}^2+
                              2c_-(p)\|q_B\|_{\rm HS}^2.
\end{equation}
The fundamental vector of the compact action is
\begin{equation}\label{eq:fundamental-action}
 \rho_z(\mathrm iS)
   =\mathrm i(S_+A-AS_-,\,S_-B-BS_+),\qquad S\in\mathfrak h.
\end{equation}
The total scalar action is ineffective, so $\mathrm i\mathfrak h$
contains all the effective compact action directions.

\begin{theorem}[Applications on the selected family]
\label{thm:integrated}
Fix the selected coefficient $h_*$, the evaluated family
of Theorem~\ref{thm:evaluated-periods}, and a finite $m\ge4$.
On the common domain of Proposition~\ref{prop:analytic-reapplication}
the following conclusions hold.

\begin{enumerate}
\item For every $p\in B_{\rm an}^{\circ}$, the integral map $D$
recovered by Theorem~\ref{thm:integral-reconstruction} is the
vanishing-cycle map in the attachment sequence, and
\begin{equation}\label{eq:attachment-hodge-identity}
 \im\mathsf r=\ker(D\otimes\RR),\qquad
 \ker D=\ZZ v_f\oplus\ZZ v_g,\qquad
 E_H(p)(x)=\sum_i\kappa_i(p)|(\mathsf r x)_i|^2.
\end{equation}
With mass $M_H$ in \eqref{eq:hodge-model}, this energy controls
the complete positive weak two-form spectrum and its observed
forms as in Recollection~\ref{thm:hodge-input}, for the actual
Ricci-flat metric on this $X_p$.

\item On the same $X_p$, the rank-four bundles in
Recollection~\ref{thm:bundle-input} have $c_2H=32$ and the stated
semistable and polystable types. For every sufficiently small
$z\in\mu_{2,p}^{-1}(0)$ put $w=\chi_p(z)$ and fix its actual
HE metric $h_0$. The invariant graph of
Recollection~\ref{thm:flow-input} defines
\[
 I(S)=h_0e^{Z(S)},\qquad Z(S)=S+\psi(S).
\]
On its actual HYM region, its finite field is the negative
gradient of twice the Donaldson functional restricted to $I$,
for the pullback of the invariant $L^2$ metric. At the HE origin,
in logarithmic metric coordinates,
\begin{equation}\label{eq:exact-spectral-flow}
 \begin{split}
 J_0&=I_{\mathfrak h}+D\psi(0)
                  =(\Pi|_{S_{B,p,w}})^{-1},\\
 J_0Df(0)&=-\mathcal L_{p,w}J_0,\\
 g_{{\rm fin},0}(S,T)
      &=M_B(S,T)+\langle D\psi(0)S,D\psi(0)T\rangle_2 .
 \end{split}
\end{equation}
Here $S_{B,p,w}$ is the entire seven-dimensional low section
space, including its kernel. If $H_{\rm fin}$ is this finite
Hessian, then
\begin{equation}\label{eq:integrated-flow-hessian}
 |g_{{\rm fin},0}-M_B|\le C|z|^4M_B,\qquad
 (1-C|z|)E_B\le H_{\rm fin}\le(1+C|z|)E_B.
\end{equation}
In particular its kernel is exactly $\ker E_B$, and
\[
 |(H_{\rm fin}-E_B)(S,T)|
       \le C|z|\sqrt{E_B(S)E_B(T)}.
\]
The complete low spectrum in Recollection~\ref{thm:bundle-input}
and the all-time actual-flow comparison in
Recollection~\ref{thm:flow-input} consequently use these same
finite metric coordinates.

\item The weak bundle energy is the norm of the action map:
\begin{equation}\label{eq:action-energy-identity}
 E_B(p,z)(S)=G_{0,p}(\rho_z(\mathrm iS),\rho_z(\mathrm iS)).
\end{equation}
For genuine tangents $q$ to a zero-level orbit-type stratum,
the actual harmonic connection metric is relatively
$1+O(|z|)$ equivalent to
\begin{equation}\label{eq:integrated-quotient}
 \inf_{\xi\in\mathrm i\mathfrak h}
       G_{0,p}(q-\rho_z\xi,q-\rho_z\xi).
\end{equation}
Thus $\ker\rho_z$ describes the constant infinitesimal
stabilizer, whereas $\im\rho_z$ describes the vertical
deformation tangents. The genuine class map, confined
distances, completion and tangent cone are those of
Recollection~\ref{thm:moduli-input}.

\end{enumerate}

The comparisons are uniform over arbitrary node-capacity
ratios, arbitrary positive bundle-rate ratios and independent
base and bundle limits in the stated domain. Observation
paths may lie in the larger $B_{\rm ev}^{\circ}$. The analytic
radius need not be computable, and no uniformity over varying
central coefficients is asserted.
\end{theorem}

The Hodge, bundle, flow and harmonic metric estimates are
imported in their full stated forms from \cite{W_Hodge,W_HYM},
with their hypotheses checked in
Proposition~\ref{prop:analytic-reapplication}. The proof below
establishes their compatibility with the integral construction.
It also records the spectral--flow and action-map
identifications explicitly; the corresponding analytic
arguments are in \cite[Sections~10--12]{W_HYM}.

\subsection{The actual attachment sequence}

The integral Thom sequence \eqref{eq:integral-attachment}
identifies the restriction of a divisor class to the $i$th
link with its degree on the $i$th exceptional curve.
The next map sends that link generator to the Thom class
Poincar\'e dual to the oriented vanishing sphere. After the
duality and excision identifications, it is exactly
$e_i\mapsto\delta_i$. Its matrix in the geometric marking
constructed here is $D$, because
Theorem~\ref{thm:integral-reconstruction} recovers the
coordinates of these spheres from their actual period jumps.

The image of the integral restriction map is
$\ZZ v_f\oplus\ZZ v_g$, as computed in
Section~\ref{sec:algebraic-family}. Exactness gives
the kernel assertion in \eqref{eq:attachment-hodge-identity}.
Over $\RR$, the kernel of restriction is the full surviving
two-dimensional space $K_H$. Its quotient therefore gives
\eqref{eq:lost-restriction} and $\im\mathsf r=\ker(D\otimes\RR)$.
The normalized link generator has capacity $\kappa_i$.
The sum of its squared restrictions weighted by these
capacities is precisely $E_H$. Minimizing the central mass
over all representatives modulo $K_H$ gives $M_H$.
The actual-metric comparison then follows from
Recollection~\ref{thm:hodge-input}, including its complete
kernel and complementary gap.

In particular, the two lost degree-two directions and the
six-dimensional vanishing-cycle span are distinct terms
of this exact sequence. The recovered placement of that
span in $\ZZ^{134}$ and its individual columns add marked
information beyond the known kernel of $D$. Its primitivity
and unmarked normal form are already determined by topology.
No small-spectrum measurement has been used to perform
this integer recovery.

\subsection{The graph tangent and the finite Hessian}

For the space of determinant-normalized Hermitian metrics
on the fixed holomorphic bundle, let
\[
 \mathcal G_h(a,b)=\int_{X_p}
       \tr(h^{-1}a\,h^{-1}b)\,\dd V_{g_p}.
\]
The HYM field $h^{-1}\dot h=-2\mathrm i\Lambda F_h$ is
the negative gradient of $\mathcal F=2\mathcal D$, where
$\mathcal D$ is Donaldson's functional with first variation
$\dd\mathcal D_h(a)=\int\tr(\mathrm i\Lambda F_h\,h^{-1}a)\dd V$.
On the inner region of the invariant graph,
$dI\,f=-\operatorname{grad}_{\mathcal G}\mathcal F\circ I$.
For every finite tangent $a$ this implies
\[
 (I^*\mathcal G)(f,a)
   =\mathcal G(dI\,f,dI\,a)
   =-\dd\mathcal F(dI\,a)
   =-\dd(\mathcal F\circ I)(a).
\]
This proves the finite gradient assertion on its actual
domain; no such assertion is needed for the auxiliary
cutoff field outside that domain.

The graph construction gives $\Pi D\psi=0$ and
$D\psi=O(|w|^2)$ in $X$. Its linearized invariance equation
at the critical origin gives
$J_0Df(0)=-\mathcal L_{p,w}J_0$.
The image of $J_0$ is in the elliptic operator domain and
is invariant under the self-adjoint $\mathcal L_{p,w}$.
The bound $Df(0)=O(|w|^2)$ puts its seven eigenvalues
below the fixed complementary gap. It is therefore
the whole low space. Since $\Pi J_0=I_{\mathfrak h}$,
$J_0$ is exactly the inverse restricted observer in
\eqref{eq:exact-spectral-flow}.

At a critical point the second derivative of $I$ is
paired with $\dd\mathcal F=0$ in the Hessian chain rule.
Consequently the exact observed forms are
\begin{equation}\label{eq:observed-graph-forms}
 N(S,T)=\langle J_0S,J_0T\rangle_2,\qquad
 H_{\rm fin}(S,T)=
              \langle\mathcal L_{p,w}J_0S,J_0T\rangle_2.
\end{equation}
Orthogonality of $\mathfrak h$ and $QX$ proves the mass
identity and its relative $O(|w|^4)$ bound.

For completeness, the passage from ordered eigenvalues
to the entire Hessian retains all weak directions.
The primitive section estimate of
\cite[Section~6]{W_HYM} is, with $r=|w|$,
\begin{equation}\label{eq:primitive-bundle-form}
 (1-Cr)\{E_B(p,w)(S)+q_{p,0}(v)\}
 \le q_{p,w}(S+v)
 \le(1+Cr)\{E_B(p,w)(S)+q_{p,0}(v)\},
\end{equation}
for every $v\perp\mathfrak h$ in $H^1$, where
$q_{p,w}(u)=\|D_{p,w}u\|_2^2$ and
$q_{p,0}(v)\ge c\|v\|_2^2$. The energy notation
$E_B(p,w)$ means the formula in \eqref{eq:bundle-model}
with arrows $w_A,w_B$.
Taking $S+v=J_0S$ gives
$H_{\rm fin}\ge(1-Cr)E_B(p,w)$.
The kernels agree: both correspond to the constant
Hermitian stabilizer, and $J_0$ fixes those constant
kernel directions.

Pass to the quotient by this kernel and denote the
induced energies by $\bar H,\bar E$ and the quotient
masses by $\bar N,\bar M$. A quotient mass is the
infimum over representatives. It preserves the comparison
$N=(1+O(r^4))M_B$.
The positive generalized eigenvalues of $(H_{\rm fin},N)$
are the positive low eigenvalues of $\mathcal L_{p,w}$.
Their relative comparison in the $w$ coordinates,
proved from \eqref{eq:primitive-bundle-form}, gives
\[
 \frac{\det\bar H}{\det\bar E}
  =\frac{\prod_{\lambda_j>0}\lambda_j}
          {\prod_{\nu_j>0}\nu_j}
       \frac{\det\bar N}{\det\bar M}
  =1+O(r).
\]
Every eigenvalue of $\bar E^{-1/2}\bar H\bar E^{-1/2}$
is at least $1-Cr$, their product is $1+O(r)$,
and their number is at most seven. Each is therefore
at most $1+Cr$. For a zero-dimensional quotient the
claim is immediate. This gives the whole relative
Hessian bound without a lower bound on a positive
weak rate.

Finally real $K_B$-equivariance implies
\[
 \rho_w(\mathrm iS)=d\chi_p(z)\rho_z(\mathrm iS).
\]
The derivative is $I+O(|z|)$ in the uniformly positive
norm $G_{0,p}$. Taking its squared norm proves
$E_B(p,w)=(1+O(|z|))E_B(p,z)$ as quadratic forms,
with the same kernel. This proves
\eqref{eq:integrated-flow-hessian}. On the quotient by that
kernel the normalized symmetric error has operator
norm at most $C|z|$; Cauchy--Schwarz proves the stated
bilinear estimate.

The phase relation is consistent with this identification.
After the basin shrink in Recollection~\ref{thm:flow-input},
both the solution shifted by $t$ and its tracking graph
trajectory remain in the phase domain. The exponential
tracking phase is unique: two finite trajectories tracking
at rate $c$ have initial separation bounded by their
separation at time $T$ times $e^{L_fT}$, which tends to
zero since $L_f<c$. Thus
$\Gamma(u(t))=\phi_t\Gamma(u(0))$.
On graph data $\Gamma(Z(S))=S$.
This identifies the actual evolution and its full-data
phase; it does not replace that phase by the initial
projection $\Pi u_0$.

\subsection{The shared action and the harmonic metric}

Equation~\eqref{eq:action-energy-identity} follows directly
from \eqref{eq:bundle-model} and
\eqref{eq:fundamental-action}. Its metric meaning requires
minimizing over all spatial gauge directions.
For the actual connection map $\mathcal A_p(w)$, the
estimate proved in \cite[Section~12]{W_HYM} is
\begin{equation}\label{eq:spatial-gauge-comparison}
 \begin{split}
 (1-C|w|)\{G_{0,p}(q,q)+\|D_0v\|_2^2\}
 &\le\|d\mathcal A_p(w)q-D_{\mathcal A_p(w)}v\|_2^2\\
 &\le(1+C|w|)\{G_{0,p}(q,q)+\|D_0v\|_2^2\},
 \end{split}
\end{equation}
for every $q\in\mathcal W$ and complementary trace-free
skew-Hermitian $v\in QH^1$.
Its proof uses harmonic orthogonality of the linear
connection increment to $D_0v$, the connection remainder
of order $O(|w|)$ relative to these two terms, and the
uniform complementary section gap.

Infimizing first in $v$ preserves the lower bound by
$G_{0,p}$; taking $v=0$ proves the upper bound.
Every remaining gauge direction is a constant
$\xi\in\mathrm i\mathfrak h$. Exact equivariance gives
$d\mathcal A_p(w)(\rho_w\xi)=-D_{\mathcal A_p(w)}\xi$.
Hence, after reversing the sign of $\xi$, apply the
same inequality to every residual $q-\rho_w\xi$ and
then take its infimum. The finite image of $\rho_w$
is closed at each fixed parameter. No bound for the
norm of a minimizing $\xi$, and no inverse of a small
positive orbit singular value, is required.

The identity $d\chi_p(z)\rho_z\xi=\rho_w\xi$ transfers
this comparison to $z$ inside each residual, before
minimization. For a tangent to an actual HE orbit-type
stratum, the full minimum is the Coulomb representative.
The linearized integrability and mean-curvature equations,
together with the K\"ahler identities, make its $(0,1)$
part Dolbeault harmonic. Its real connection norm is twice
the Dolbeault harmonic norm. This is exactly the convention
used in \eqref{eq:arrow-mass} and
Recollection~\ref{thm:moduli-input}.
The actual class and stabilizer assertions in that theorem
then give the genuine quotient interpretation.
This proves \eqref{eq:integrated-quotient} through every
stabilizer type.

In particular the flow graph metric acts on metrics
of one fixed holomorphic bundle, while the harmonic
metric acts on deformation tangents modulo $\im\rho_z$.
Their relationship is the common action map and its
stabilizer, rather than equality of the two metrics.

\subsection{Independence of the integral observations from the bundle metric}

For fixed $p,w$, unitary bundle gauge transformations and
HYM metric evolution leave the equations of $X_p$, the
residue form, the marked cycles and their transports unchanged.
The observations $\Delta_{ai}=p_i(Jd_i)_a$ and the recovered
map $D$ are therefore independent of those bundle operations.
They use the represented cycles and period enclosures of
Theorem~\ref{thm:accessible-observations}. Changes of homology
marking and signed node labels instead obey the covariance
and error transformations of Section~\ref{sec:integral-recovery}.

The same node labels enter $\mathsf r$, the periods
$p_i$, and the smoothing functions $\tau_i=p_i/U_i$.
Thus \eqref{eq:capacity-digits} relates the required
absolute period precision to the actual capacity
weights. It does not identify period error with the
relative spectral error $e(p)$. Proposition
\ref{prop:analytic-reapplication} ensures that both
constructions concern these same target fibres even
when observation paths leave $B_{\rm an}$.

\subsection{Independent base and bundle scales}

\begin{corollary}\label{cor:independent-scales}
Let $F(h),G(h)\to+\infty$ and set
\[
 p(h)=(a,a,a,-3a;\ b,b,b,-3b),\qquad
 a=e^{-F(h)},\quad b=e^{-G(h)}.
\]
For all sufficiently small $h$ these are actual smooth
fibres of the selected family, with
$w_f\sim8\pi/F$ and $w_g\sim8\pi/G$.
If $F/G\to0$, their larger and smaller positive weak
Hodge rates satisfy respectively
\[
 \lambda_+\sim\frac{9\pi}{8F},\qquad
 \lambda_-\sim\frac{10\pi}{9G}.
\]
Independently, for any $q>1$ and $s\downarrow0$, take
\[
 \sqrt{c_+(p)}\,A=\sqrt{c_-(p)}\,B
                              =\diag(s,s+s^q).
\]
Their actual HE images have one trace-free zero mode,
two positive rates
$(2/V_H)s^{2q}(1+O(s))$, and four rates asymptotic to
$(8/V_H)s^2$. The actual HYM comparison holds through
times of order $s^{-2q}$ and all later times.
The harmonic quotient and integral reconstruction
conclusions apply on these same fibres.
\end{corollary}

\begin{proof}
The two sum-zero relations hold identically and no
coordinate vanishes. The evaluated full period ball
contains all sufficiently small points of this path.
The residue units are bounded above and below, so
the logarithms in each group differ from $F$ or $G$
by bounded quantities. This proves the capacity
asymptotics. Formula~\eqref{eq:hodge-roots}, or its
trace and determinant, gives the two displayed
Hodge rates when $F/G\to0$; relative actual-metric
comparison preserves them.

The weighted arrows are equal diagonal matrices,
so their quadratic moment map vanishes exactly.
For positive diagonal entries $u_1,u_2$, the seven
finite section rates are
\[
 0,\quad \frac{8u_1^2}{V_H},\quad\frac{8u_2^2}{V_H},
 \quad\frac{2(u_1-u_2)^2}{V_H}\ \text{(twice)},\quad
 \frac{2(u_1+u_2)^2}{V_H}\ \text{(twice)}.
\]
This follows by separating diagonal Hermitian entries
and the real and imaginary parts of each off-diagonal
entry in \eqref{eq:bundle-model}; the weighted arrow
masses cancel. Substitution gives the asserted
hierarchy, and Recollection~\ref{thm:bundle-input} gives
the relative $O(s)$ actual comparison.
The domains of Recollection~\ref{thm:flow-input} and
Theorem~\ref{thm:integrated} impose no relation between
$h$ and $s$, proving the remaining assertions.
\end{proof}

Finite-order holomorphic arcs have a single logarithmic
order, with coefficients depending on their vanishing
orders. The paths in the corollary may have genuinely
separated logarithmic scales and need not be holomorphic
at the origin. When the ratio of the group scales has
no limit, the pointwise comparisons remain valid;
limiting hierarchies require subsequences or parameter
sectors. The restriction and action maps above do not
by themselves construct a filtered elliptic complex.

The integrated theorem is on smooth base fibres,
polystable bundle parameters and the stated spatial
metric basin. Singular graded concentration requires
additional local data and is a separate problem.
The theorem supplies neither singular-fibre HYM
evolution nor a completion of a larger moduli space.
Its finite evolution is the exact corrected field,
and its all-time error is absolute; a relative
nonlinear decay expansion requires further estimates.
